\subsection{\texorpdfstring{特征为 \(p\) 的完全环}{特征为 p 的完全环}}

定义 2. --- 一个特征为 \(p \neq 0\) 的环 A 称为完全环，如果它是交换的且映射 \(a \mapsto a^p\) 是双射的。

如果环 A 是特征 p 的完全环，映射 \(a \mapsto a^{p^{f}}\) 对每个整数 \(f \geqslant 0\) 是环 A 的自同构；逆自同构记为 \(a \mapsto a^{1/p^{f}}\) 或 \(a \mapsto a^{p^{-f}}\)，并且 A 的子集 S 在该自同构下的像记为 \(S^{1/p^{f}}\) 或 \(S^{p^{-f}}\) 。显然 \((a^{p^{e}})^{p^{f}} = a^{p^{e+f}}\) 对所有 \(a \in A\) 以及任意整数 e 和 f（无论其符号如何）成立。

设 A 是特征为 p 的交换环。对每个整数 \(f \geqslant 0\)，用 \(n_{f}\) 表示环 A 的自同态 \(a \mapsto a^{p^{f}}\) 的核。则 \((\mathfrak{n}_{f})_{f \geqslant 0}\) 是 A 的理想的一个递增序列；由于每个正整数都被 p 的某个幂所控制，理想 \(n = \bigcup_{f \geqslant 0} n_{f}\) 由 A 的所有

幂零元素组成。特别地，若 A 是完全的，则 A 的每个幂零元素都为零。

定义 3. --- 设 A 是一个特征为 \(p \neq 0\) 的交换环。所谓 A 的完全闭包，我们指的是一个对 \((\hat{\mathbf{A}}, u)\)，其中 \(\hat{\mathbf{A}}\) 是特征为 \(p\) 的完全环，且 \(u\) 是 A 到 \(\hat{\mathbf{A}}\) 中的一个同态，满足如下泛性质：

（PC）若 B 是特征为 p 的完全环，且 v 是从 A 到 B 的同态，则存在唯一的同态 h（从 \(\hat{A}\) 到 B 中），使得 \(v = h \circ u\) 。

泛性质（PC）立即蕴含完全闭包的唯一性，其含义如下：若 \((\hat{\mathbf{A}}, u)\) 与 \((\hat{\mathbf{A}}', u')\) 是 A 的两个完全闭包，则存在唯一的、把 \(\hat{A}\) 映到 \(\hat{A}'\) 上的同构 h，使得 \(u' = h \circ u\) （参见 E，IV，第 23 页）。我们现在来建立完全闭包的存在性：

定理 3. --- 设 A 是一个特征为 \(p \neq 0\) 的交换环。则 A 存在一个完全闭包 \((\hat{\mathbf{A}}, u)\)。此外，\(u\) 的核是 A 的所有幂零元素组成的集合，并且对于每个 \(x \in \hat{\mathbf{A}}\)，存在一个整数 \(n \geqslant 0\)，使得 \(x^{p^n} \in u(\mathbf{A})\) 。

对每个整数 \(n \geqslant 0\)，令 \(\mathbf{A}_n = \mathbf{A}\)；当 \(m \geqslant n\) 时，我们定义一个同态 \(\pi_{m,n}\)（从 \(\mathbf{A}_n\) 到 \(\mathbf{A}_m\) 中），由 \(\pi_{m,n}(a) = a^{p^{m-n}}\) 给出。这样我们得到环的一个正向系统 \((\mathbf{A}_n, \pi_{m,n})\)（I，第 120 页）；设 \(\hat{\mathbf{A}}\) 为该系统的正向极限，\(u_n\) 为 \(\mathbf{A}_n = \mathbf{A}\) 到 \(\hat{\mathbf{A}}\) 中的典范同态；我们还令 \(u = u_0\)。根据正向极限的构造，核 \(n\) （\(u\) 的）是 A 到 A 中的同态 \(\pi_{n,0}: a \mapsto a^{p^n}\) 的核的并集，因此它由 A 的所有幂零元素组成。环 \(\hat{A}\) 根据 V 第 3 页的注记 1 是特征为 p 的交换环。

环 \(\hat{\mathbf{A}}\) 是子环的递增序列 \((u_n(\mathbf{A}))_{n\geq 0}\) 的并集。我们有 \(u_{n}(\mathbf{A})^{p^{n}} = \mathbf{A}\)；因此对每个 \(x\in \hat{\mathbf{A}}\)，存在一个整数 \(n\geqslant 0\)，使得 \(x^{p^n}\in u(\mathbf{A})\)。我们还有 \(u_{n}(\mathbf{A}) = u_{n + 1}(\mathbf{A})^{p}\)，从而 \(\hat{\mathbf{A}}^p = \hat{\mathbf{A}}\)。设 \(x\in \hat{\mathbf{A}}\) 使得 \(x^{p} = 0\)；选择一个整数 \(n\geqslant 1\) 和一个元素 \(a\in \mathbf{A}\)，使得 \(x = u_{n}(a)\)。于是我们有 \(u_{n - 1}(a) = u_n(a)^p = 0\)；根据正向极限的定义，存在一个整数 \(m\geqslant n\)，使得 \(\pi_{m - 1,n - 1}(a) = 0\)，即 \(a^{p^{m - n}} = 0\)。因此我们有 \(\pi_{m,n}(a) = 0\)，从而 \(u_{n}(a) = 0\)，即 \(x = 0\)。因此环 \(\hat{\mathbf{A}}\) 是特征为 \(p\) 的完全环。

设 \(v\) 是 \(A\) 到完全环 \(B\)（特征为 \(p\)）中的一个同态。对每个整数 \(n \geqslant 0\)，映射 \(b \mapsto b^{p^n}\) 是 \(B\) 的一个自同构，因此存在一个同态 \(v_n\) （从 \(A_n = A\) 到 \(B\) 中），它由 \(v(a) = v_n(a)^{p^n}\) 刻画。于是我们有 \(v_m \circ \pi_{m,n} = v_n\)，对 \(m \geqslant n \geqslant 0\)；根据正向极限的定义，存在一个同态 \(h\) （从 \(\hat{A}\) 到 \(B\) 中），使得 \(v_n = h \circ u_n\) 对所有 \(n \geqslant 0\) 成立；特别地，我们有 \(v = v_0 = h \circ u_0 = h \circ u\)。最后设 \(h'\) 是 \(\hat{A}\) 到 \(B\) 中的一个同态，使得 \(h' \circ u = v\)。设 \(x \in \hat{A}\)；如我们所见，存在一个整数 \(n \geqslant 0\) 和一个元素 \(a \in A\)，使得 \(x^{p^n} = u(a)\)。于是我们有

\[
h (x) ^ {p ^ {n}} = h (u (a)) = v (a) = h ^ {\prime} (u (a)) = h ^ {\prime} (x) ^ {p ^ {n}},
\]

且由于 B 是完全的，我们得到 \(h(x) = h'(x)\)。于是我们有 \(h' = h\)，这就完成了 \((\hat{\mathbf{A}}, u)\) 是 A 的完全闭包的证明。

命题 3. --- 设 B 是特征为 p 的完全环，A 是 B 的一个子环。记 \(A^{p^{-\infty}} = \bigcup_{f \geqslant 0} A^{p^{-f}}\)，并用 j 表示 A 到

\(\mathbf{A}^{p^{-\infty}}\) 中的典范单射。则 \(\mathbf{A}^{p^{-\infty}}\) 是 \(\mathbf{B}\) 的包含 \(\mathbf{A}\) 的最小完全子环，且 \((\mathbf{A}^{p^{-\infty}}, j)\) 是 \(\mathbf{A}\) 的一个完全闭包。

对每个整数 \(f \in \mathbf{Z}\)，用 \(\pi_f\) 表示自同构 \(b \mapsto b^{p^f}\) （\(\mathbf{B}\) 的）。子环序列 \(\pi_{-f}(\mathbf{A})\) （在 \(\mathbf{B}\) 中）（对 \(f \geqslant 0\)）是递增的，其并集 \(\mathbf{A}^{p^{-\infty}}\) 因而是 \(\mathbf{B}\) 的一个子环。我们有 \(\pi_1(\mathbf{A}^{p^{-\infty}}) = \bigcup_{f \geqslant 0} \pi_{-(f-1)}(\mathbf{A}) = \mathbf{A}^{p^{-\infty}}\)，因此

\(A^{p^{-\infty}}\) 是 B 的一个完全子环。最后令 \(B_0\) 是 B 的一个完全子环且包含 A；对每个整数 \(f \geqslant 0\)，我们有 \(\pi_{-f}(A) \subset \pi_{-f}(B_0) = B_0\)，从而 \(A^{p^{-\infty}} \subset B_0\) 。

如果 \(v\) 是 \(\mathbf{A}\) 到完全环 \(\mathbf{B}'\)（特征为 \(p\)）中的一个同态，则对每个整数 \(f \geqslant 0\)，我们可以定义一个同态 \(h_f\) （从 \(\pi_{-f}(\mathbf{A})\) 到 \(\mathbf{B}'\) 中），由 \(h_f(\pi_{-f}(a)) = v(a)^{p^{-f}}\)（对所有 \(a \in \mathbf{A}\)）给出。我们立刻看到 \(h_{f + 1}\) 与 \(h_f\) 在 \(\pi_{-f}(\mathbf{A})\) 上一致；因此存在一个同态 \(h\) （从 \(\mathbf{A}^{p^{-\infty}}\) 到 \(\mathbf{B}'\) 中），它把 \(h_f\) 诱导在 \(\pi_{-f}(\mathbf{A})\) 上（对所有 \(f \geqslant 0\)），特别地，\(h\) 扩张 \(h_0 = v\)。若 \(h'\) 是 v 的另一个扩张——扩张为 \(A^{p^{-\infty}}\) 到 B 中的同态——则等式 \(h' = h\) 可按定理 3 的证明中的方法建立。

